\section{Step 23: Publish the accepted frame}
Only an accepted, fully applied image may advance the decoder's logical
state. Read the saved expected sequence, add one, and store it in context offset
eight. Store the accepted kind at offset 24 and the active image byte count,
16\,384, at offset 32. Set the success return values. The supplied epilogue
restores the caller's saved registers and stack pointer.

Fill \texttt{record\_publish}. Keep capacity, period and the suffix unchanged.
Run full, delta, unchanged, optional-full and invalid controls. Then call the
routine repeatedly in one execution session. Insert an invalid late record
between valid records: it must retain the last accepted frame and the next
expected sequence, allowing a subsequent valid record with that same sequence
to succeed. The record interface now has useful behavior independently of the
future file reader, directory and playback loop.

\paragraph{Observe} Prepare \texttt{s23-1} and stop before
\texttt{record\_published}. At \texttt{record\_published}, \texttt{\$v0} is zero and
\texttt{\$v1} is 16\,384. The worked context now has next sequence two,
last kind one and active bytes 16\,384. The caller prints
\texttt{Record applied.} after returning.

\Needspace{7\baselineskip}
\paragraph{Check} What happens to the expected sequence on failure?
\begin{itemize}
\item[A.] It advances to avoid retrying the bad record.
\item[B.] It remains unchanged.
\item[C.] It becomes the record payload length.
\end{itemize}
